\section{Circuit depth, lower bounds and barriers}

Circuit classes in this section are nonuniform. Unless a different cost is
specified, native De Morgan circuits charge constants, NOT, AND and OR gates,
with wiring and fanout free. Formula leaf count and formula depth are separate
measures; unbounded fan-in bases are identified in their entries.

\begin{result}
\subsection{Shannon--Lupanov bounds and the MCSP window}
Count all NOT, AND, OR and constant gates. For every $\varepsilon>0$, all
sufficiently large $n$, and every $n$-input truth table, MCSP answers yes at
any integer threshold $s\ge(1+\varepsilon)2^n/n$. For every sufficiently large
$n$, some truth table needs more than $\lfloor2^n/n\rfloor$ gates.
A finite counting bound makes the scarcity of small circuits explicit: for
$n\ge16$, at most $2^{\lfloor2^n/4\rfloor}$ of the $2^{2^n}$ truth tables
have a circuit of size at most $\lfloor2^n/(5n)\rfloor$.
These are synthesis and counting statements; they do not establish the
computational hardness of deciding MCSP.\citep{Shannon1949,Lupanov1958}
\decl{Complexity/Computability/MCSP/Shannon.lean}{248}{Cslib.MCSP.eventually_mem}
\decl{Complexity/Computability/MCSP/Shannon.lean}{256}{Cslib.MCSP.eventually_exists_not_mem}
\decl{Complexity/Computability/MCSP/Shannon.lean}{226}{Cslib.MCSP.ncard_setOf_mem_le}
\end{result}

\begin{result}
\subsection{Barrington's theorem}
Nonuniform NC$^1$ is exactly the class of languages recognized by polynomial-length
permutation branching programs of width five. An equivalent formulation uses
programs over the alternating group $A_5$. This is a class equality with the
library's nonuniform circuit families; it does not assert a uniform-family
characterization.\citep{Barrington1989,BarringtonTherien1988}
\decl{Complexity/Computability/Circuit/Boolean/BarringtonClass.lean}{205}{Cslib.Circuits.Boolean.NC_one_eq_polyLength_perm_five}
\decl{Complexity/Computability/Circuit/Boolean/BarringtonClass.lean}{217}{Cslib.Circuits.Boolean.NC_one_eq_polyLength_alternatingGroup}
\end{result}

\begin{result}
\subsection{Spira's formula balancing theorem}
Every De Morgan formula with $L\ge1$ input leaves has an equivalent formula with
at most $4\lceil\log_2 L\rceil$ AND/OR gates on each input-to-output path.
Constants contribute zero leaves; when $L=0$, the formal bound is zero. Negations have
zero cost in this depth measure. The theorem certifies the depth of the balanced
formula, not a machine runtime for constructing it.\citep{Spira1971}
\decl{Complexity/Computability/Circuit/Boolean/Spira.lean}{211}{Cslib.Circuits.Boolean.exists_formula_depthWith_le_four_mul_clog}
\end{result}

\begin{result}
\subsection{Khrapchenko's bound and exact parity formula size}
Let a De Morgan formula be true on $A$ and false on $B$, and let $E(A,B)$ be the
pairs in $A\times B$ at Hamming distance one. Its input-leaf count $L$ (constants contribute zero) satisfies
\[
|E(A,B)|^2\le |A|\,|B|\,L.
\]
In particular, parity on $n$ bits requires at least $n^2$ leaves; for $n=2^k$
this is attained. Yablonskii's construction gives, for $n=2^k+r$ with
$0\le r\le2^k$, a parity formula with at most $2^k(2^k+3r)$ leaves;
for every $n$ there is one with $8L\le9n^2$.
These are formula leaf bounds, not bounds on circuits with shared subcomputations.
\citep{Khrapchenko1971b,Khrapchenko1971,Yablonskii1954,Jukna2012}
\decl{Complexity/Computability/Circuit/Boolean/Khrapchenko.lean}{144}{Cslib.Circuits.Formula.sq_card_hammingEdges_le}
\decl{Complexity/Computability/Circuit/Boolean/Khrapchenko.lean}{428}{Cslib.Circuits.Boolean.isLeast_leafSize_parity_two_pow}
\decl{Complexity/Computability/Circuit/Boolean/Yablonskii.lean}{167}{Cslib.Circuits.Boolean.exists_formula_parity_two_pow_add}
\decl{Complexity/Computability/Circuit/Boolean/Yablonskii.lean}{182}{Cslib.Circuits.Boolean.exists_formula_parity_eight_mul_leafSize_le}
\end{result}

\begin{result}
\subsection{Parity lower bounds by random restrictions}
A depth-$d$ circuit of unbounded fan-in AND/OR gates and negations computing
parity on $n$ bits has at least $2^s$ gates whenever
\[
s\ge2,\qquad s(16s+1)^d\le n.
\]
Thus the finite theorem gives an exponential lower bound at every fixed depth.
Depth counts negations and size counts every gate. The exponent obtained is of
order $n^{1/(d+1)}$, weaker than the sharp classical H\aa stad exponent
$n^{1/(d-1)}$, which the next entry proves.\citep{FurstSaxeSipser1984,Ajtai1983,Hastad1986}
\decl{Complexity/Computability/Circuit/Boolean/Hastad.lean}{367}{Cslib.Circuits.Boolean.Unbounded.two_pow_le_size_of_computes_parity}
\end{result}

\begin{result}
\subsection{H\aa stad's theorem with the sharp exponent}
Let $C$ be a circuit of negations and unbounded fan-in AND/OR gates of AND/OR depth at most $d$
(the largest number of AND and OR gates on a path; negations are free, and constants are nullary
AND/OR gates) computing parity on $n$ bits. Then $C$ has at least $2^s$ AND/OR gates whenever
\[ s\ge2,\qquad 17s(16s+1)^{d-2}\le n \]
($d-2$ truncated at $0$). Hence for $d\ge2$ it has at least $2^s$ gates whenever
$(17s)^{d-1}\le n$, and at least $\tfrac12\,2^{n^{1/(d-1)}/17}$ gates once $n\ge34^{d-1}$.
The first bound also holds when depth counts negations, and for formulas. With the
Lecomte--Ramakrishnan upper bound $2^{O(n^{1/(d-1)})}+n$ already in the library, the exponent
$1/(d-1)$ is optimal. The proof is the staged switching argument with stages indexed by AND/OR
depth, a first stage switching $1$-DNFs (keeping a $1/17$ fraction of the variables), and no final
stage: a DNF whose terms are shorter than the number of free variables cannot compute restricted
parity.\citep{Hastad1986,Jukna2012,AroraBarak09,LecomteRamakrishnan2026}
\declat{fa8e00652d0d88d6792118a996e01f1031cbdc4f}{Complexity/Computability/Circuit/Boolean/Hastad/Sharp.lean}{541}{Cslib.Circuits.Boolean.Unbounded.two_pow_le_card_andOr_of_computes_parity}
\declat{fa8e00652d0d88d6792118a996e01f1031cbdc4f}{Complexity/Computability/Circuit/Boolean/Hastad/Sharp.lean}{551}{Cslib.Circuits.Boolean.Unbounded.two_pow_le_size_of_computes_parity_of_depthWith_le}
\declat{fa8e00652d0d88d6792118a996e01f1031cbdc4f}{Complexity/Computability/Circuit/Boolean/Hastad/Sharp.lean}{561}{Cslib.Circuits.Boolean.Unbounded.two_pow_le_size_of_computes_parity_of_pow_le}
\declat{fa8e00652d0d88d6792118a996e01f1031cbdc4f}{Complexity/Computability/Circuit/Boolean/Hastad/Sharp.lean}{572}{Cslib.Circuits.Boolean.Unbounded.two_rpow_le_two_mul_size_of_computes_parity}
\declat{fa8e00652d0d88d6792118a996e01f1031cbdc4f}{Complexity/Computability/Circuit/Boolean/Hastad/Sharp.lean}{608}{Cslib.Circuits.Boolean.Unbounded.two_pow_le_size_of_computes_parity_of_depth_le}
\declat{fa8e00652d0d88d6792118a996e01f1031cbdc4f}{Complexity/Computability/Circuit/Boolean/Hastad/Sharp.lean}{617}{Cslib.Circuits.Boolean.Unbounded.two_pow_le_size_of_formula_computes_parity_of_depthWith_le}
\end{result}

\begin{result}
\subsection{Parity is not first-order definable}
Let $L$ be a relational vocabulary with finitely many relation symbols, $P\in L$ unary, and
$N_n$ an arbitrary $L$-structure on $\{0,\dots,n-1\}$ for each $n$ (numerical relations such as
order or BIT, chosen separately per size). Write $N_n[P{:=}S]$ for $N_n$ with $P$ reinterpreted
as $S$. No first-order sentence $\varphi$ of $L$ satisfies, for all large $n$ and all
$S\subseteq\{0,\dots,n-1\}$, $N_n[P{:=}S]\models\varphi\iff|S|$ is even; that is, parity is not
in FO[Arb]. In particular no sentence defines ``$|P|$ is even'' on all finite $L$-structures with
at least $n_0$ elements. Constant symbols are not supported. The proof substitutes input bits
and constants into the constant-depth polynomial-size formula deciding $\varphi$ and applies the
random-restriction parity bound for formulas above.\citep{FurstSaxeSipser1984,Ajtai1983,Immerman1999}
\declat{fa8e00652d0d88d6792118a996e01f1031cbdc4f}{Complexity/Computability/Circuit/Boolean/FirstOrder/Parity.lean}{274}{Cslib.Circuits.Boolean.ModelChecking.not_eventually_forall_realize_expansion_iff_even}
\declat{fa8e00652d0d88d6792118a996e01f1031cbdc4f}{Complexity/Computability/Circuit/Boolean/FirstOrder/Parity.lean}{299}{Cslib.Circuits.Boolean.ModelChecking.not_forall_realize_iff_even_card}
\end{result}


\begin{result}
\subsection{Razborov--Smolensky lower bounds}
For every odd prime $p$, parity is not in nonuniform AC$^0[p]$: adding unbounded
fan-in MOD$_p$ gates to constant-depth polynomial-size AND/OR/NOT circuits does
not suffice. Quantitatively, a depth-$d$ circuit computing parity has size at
least $2^{\ell+1}/5$ whenever $\ell\ge1$ and
\[
100\bigl(((p-1)\ell)^d+1\bigr)^2\le n+1.
\]
The prime and depth are fixed in the class separation.\citep{Razborov1987,Smolensky1987}
\decl{Complexity/Computability/Circuit/Boolean/Smolensky.lean}{60}{Cslib.Circuits.Boolean.Smolensky.two_mul_two_pow_le_size_of_computes_parity}
\decl{Complexity/Computability/Circuit/Boolean/Smolensky.lean}{89}{Cslib.Circuits.Boolean.Smolensky.parity_notMem_AC0Mod}
\end{result}

\begin{result}
\subsection{Razborov's monotone clique bound}
Let $p,w\ge2$, $w+1\le k\le n$, $p\le n$, and $2w^2\le k-1$, and put
\[
B=(p-1)^w w!,\quad
A=B^2\binom{n-w-1}{k-w-1},\quad
N=(2B+B^2)(w^2)^p(k-1)^{n-p}.
\]
Every monotone AND/OR circuit for $k$-clique on $n$ vertices has more than $s$
gates if $As<\binom nk$ and $2Ns<(k-1)^n$. This states the formalized finite
approximation-method bound with its parameters; no unrestricted Boolean-circuit
lower bound follows from it.\citep{Razborov1985}
\decl{Complexity/Computability/Circuit/Monotone/Clique/LowerBound.lean}{139}{Cslib.Circuits.Monotone.lt_size_of_computes_clique}
\end{result}

\begin{result}
\subsection{The natural-proofs barrier}
Assume a family of length-doubling generators of polynomial circuit size and a
fixed integer $d\ge1$ such that, for all sufficiently large seed lengths $k$ and
every $m$ with $m^d\le k$, no circuit of size at most $2^m$ distinguishes its
output from uniform with advantage greater than $2^{-m}$. Then no
P/poly-natural property is useful against P/poly.
Here natural means containing a property of density at least $2^{-O(n)}$ whose
truth-table language has circuits polynomial in the table length $2^n$; useful
means eventually excluding the slices of every P/poly language. The security
assumption is stronger than ordinary polynomial-time PRG security.
\citep{RazborovRudich1997,GGM1986}
\decl{Complexity/Computability/Circuit/Boolean/NaturalProof/Generator.lean}{51}{Cslib.Circuits.Boolean.not_isUseful_PPoly_of_isNatural_of_generator}
\end{result}

\begin{result}
\subsection{Exponentially small correlation with quadratic forms}
There is a nonconstructively chosen family of bilinear Boolean functions $h_n$
with the following property. Put
\[
c_* = 1+\frac{\pi}{6\arccos((1+2\sqrt2)/4)}\approx2.78125.
\]
For every real $c<c_*$, there is $\gamma>0$ such that for all sufficiently
large $n$, every Boolean circuit of fan-in at most two, over any Boolean basis,
with at most $cn$ positive-arity gates has
\[
\left|2\Pr_{x\in\{0,1\}^n}[C(x)=h_n(x)]-1\right|\le2^{-\gamma n}.
\]
The input is uniform, constants are not charged, and the cutoff and rate are
independent of the circuit. The forms are obtained by choosing submatrix-robust
matrices; no efficient construction or uniform complexity-class membership is
proved. The exact coefficient above is the library's frontier-method result;
related correlation-bound sources are credited separately.\citep{ChenKabanets2015,GKST2016,Jukna2011}
\decl{Complexity/Computability/Circuit/Frontier/Correlation.lean}{279}{Cslib.Circuits.Frontier.exists_eventually_abs_correlation_hardForm_le_of_lt}
\end{result}

\begin{result}
\subsection{Weighted-threshold activity and capacity bounds}
Consider real-weight threshold gates, and count as input-reading a gate with a
nonzero effective coefficient on an original input; repeated wires are aggregated
before testing this condition. If a circuit computes inner product modulo two
on two $m$-bit blocks, with its output taken from a gate, let $q$ be its number of
input-reading gates and $e$ the maximum number of those gates outputting true on
one input. Then
\[
2^m\le\left(\sum_{j=0}^{e}\binom qj\right)^2.
\]
The same inequality holds with $q$ counting all gates and $e$ their maximum
simultaneous activity.\citep{ThresholdCircuitRepresentations2026}
\decl{Complexity/Computability/Circuit/Boolean/WeightedThreshold/InnerProduct.lean}{75}{Cslib.Circuits.Boolean.WeightedThreshold.two_pow_le_sum_choose_sq_of_circuit_hadamardCode}
\decl{Complexity/Computability/Circuit/Boolean/WeightedThreshold/InnerProduct.lean}{99}{Cslib.Circuits.Boolean.WeightedThreshold.two_pow_le_sum_choose_gateActivity_sq_of_circuit_hadamardCode}
\end{result}

\begin{result}
A separate capacity bound applies to any fixed Boolean readout of a threshold
program's gate values computing $f$ on $n$ bits. Suppose that, across every
partition of the input coordinates, every rectangle contained in either fiber
of $f$ has a side with fewer than $K$ assignments. If $q$ gates read original
inputs in the effective sense above, then
\[
2^n<4K^2\,4^q.
\]
This statement requires the rectangle hypothesis on both fibers; it is not an
unconditional lower bound for every Boolean function.\citep{Schlesinger2026Algebraic}
\decl{Complexity/Computability/Circuit/Boolean/WeightedThreshold/CapacityCircuit.lean}{150}{Cslib.Circuits.Boolean.WeightedThreshold.two_pow_lt_mul_four_pow_inputGates}
\end{result}

\begin{result}
\subsection{Exact parity formula depth at powers of two}
Let $d(f)$ be minimum De Morgan formula depth, counting binary AND/OR gates
and making negations free. For $n\ge1$,
\[
\lceil\log_2(n^2)\rceil\le d(\mathrm{PARITY}_n)
 \le2\lceil\log_2 n\rceil,
\qquad d(\mathrm{PARITY}_{2^k})=2k.
\]
The lower bound comes from formula leaf count, and the upper bound from
the parity construction. This does not assert the exact minimum for every
non-power-of-two arity.\citep{Khrapchenko1971b,Yablonskii1954,Jukna2012}
\decl{Complexity/Computability/Circuit/Boolean/KhrapchenkoDepth.lean}{73}{Cslib.Circuits.Boolean.clog_sq_le_formulaDepth_parity}
\decl{Complexity/Computability/Circuit/Boolean/Yablonskii.lean}{194}{Cslib.Circuits.Boolean.formulaDepth_parity_le_two_mul_clog}
\decl{Complexity/Computability/Circuit/Boolean/KhrapchenkoDepth.lean}{88}{Cslib.Circuits.Boolean.formulaDepth_parity_two_pow}
\end{result}

\begin{result}
\subsection{Quadratic completeness of semifilter covers}
Let $A\subseteq U$ and let $B_1,\ldots,B_n\subseteq U$ be supplied input sets.
A semifilter on $U\setminus A$ is an upward-closed family containing the whole
complement and excluding the empty set. It is above a point $a\in A$ if it
contains every $B_i\cap(U\setminus A)$ for which $a\in B_i$.
Suppose a list of $t$ pairs of complement subsets defeats every semifilter above any point of $A$:
for some listed pair it accepts both components but not their intersection.
Then there is an acyclic set circuit computing $A$ from the supplied sets with
at most $t^2$ AND gates. Binary unions and both constants are free, and fanout
shares computed sets. Consequently the minimum AND cost is at most the square
of the minimum unrestricted pair-cover length (with value $\infty$ if no cover exists).
The construction works even for infinite $U$, empty input families, and trivial
targets: each point's $t$ pair coordinates stabilize within $t$ rounds. Selecting
free unions uses classical subset tests; no effective compiler running-time bound
is claimed. This is the native free-constant version of Cavalar--Oliveira's
quadratic completeness theorem, following the Fusion method.
\citep{CavalarOliveira2025}
\declat{dbb5eddb2bd1ce9ee9971d00f22c69c26ae0d542}{Complexity/Computability/Circuit/Fusion/Quadratic.lean}{226}{Cslib.Circuits.Fusion.PairCover.cost_toCircuit}
\declat{dbb5eddb2bd1ce9ee9971d00f22c69c26ae0d542}{Complexity/Computability/Circuit/Fusion/Quadratic.lean}{233}{Cslib.Circuits.Fusion.PairCover.eval_toCircuit}
\declat{dbb5eddb2bd1ce9ee9971d00f22c69c26ae0d542}{Complexity/Computability/Circuit/Fusion/Quadratic.lean}{258}{Cslib.Circuits.Fusion.ecomplexityFrom_le_pairCoverComplexity_sq}
\end{result}

\begin{result}
\subsection{Graph Fusion transfer and exact inequality complexity}
Let $V=\{0,1\}^n$, let $G\subseteq V\times V$, and supply every row and column
as a free source. A semifilter on $G^c$ accepts the full set, rejects the empty
set, and is upward closed. It is above an edge when it accepts the restricted
row and column containing that edge. A cover defeats every such witness with a
pair of accepted sets whose intersection is rejected; $\rho(G)$ is the minimum
number of pairs. Representing a Boolean input by two $n$-bit blocks gives
\[
\rho(G)\le \rho(f_G;\text{positive and negative literals}).
\]
The inequality also holds when both sides restrict to semi-ultrafilters (every
set or its complement is accepted), and includes trivial graphs.
For $G_{\ne}=\{(u,v):u\ne v\}$,
\[
\rho_{\mathrm{can}}(G_{\ne})=\rho(G_{\ne})=D_{\cap}(G_{\ne})=n.
\]
Canonical semifilters are the upward closures of a complementary row and
column, both nonempty. Here $D_{\cap}$ counts only binary intersections in
native acyclic circuits: unions, empty/full constants, rows, columns, wiring
and fanout are free. The equality includes $n=0$, when the graph is empty and
all three costs are zero. This is a graph intersection bound, not a total-gate
count.\citep{CavalarOliveira2025}
\declat{45d71952d39a80bf96fd8e6144d9fd9ec3ecf5c5}{Complexity/Computability/Circuit/Fusion/Graph.lean}{98}{Cslib.Circuits.Fusion.Graph.pairCoverComplexity_le_literals}
\declat{45d71952d39a80bf96fd8e6144d9fd9ec3ecf5c5}{Complexity/Computability/Circuit/Fusion/Graph/Inequality.lean}{212}{Cslib.Circuits.Fusion.Graph.inequality_complexities}
\end{result}

\begin{result}
\subsection{The fixed-distribution Fusion counting barrier}
Let $A\subseteq\{0,1\}^n$, let $v\in A$, and let $F$ be a semifilter on
$U=\{0,1\}^n\setminus A$. Suppose $F$ accepts
$\{u\in U:u_i=v_i\}$ for every coordinate $i$.
There is an explicit distribution on pairs $(S,T)$ of subsets of $U$, depending
only on $A$, which covers every such $(v,F)$ with probability at least
$1/(8n+8)$. Covering means $F(S)$ and $F(T)$ hold but $F(S\cap T)$ does not.
The sampler independently and uniformly chooses a coordinate or an extra
complement case, a Boolean polarity, and a full truth table.

More generally, for any finite family $(v_w,F_w)$ satisfying these hypotheses
and nonnegative real weights $\mu_w$, there exists a pair covering total weight
at least
\[
\frac{\sum_w\mu_w}{8n+8}.
\]
If the weights sum to one and every pair covers weight at most $d$, then
$d>0$ and $1/d\le8n+8$. This bounds the fixed-distribution counting certificate,
not minimum cover size or circuit complexity. The hitting pair can depend on
the weighted family; the sampling distribution does not. The theorem applies even
if the weights are chosen after seeing a circuit, provided the capacity bound
holds for every pair. It does not bound certificates whose capacity constraint
ranges only over that circuit's gate pairs.

The proof adapts Oliveira's exposition of Razborov's argument through successive
random truth-table cylinders. It needs only acceptance of matching literals,
not exact agreement on the opposite literals, and proves the sampling lemma
rather than assuming it. No efficient sampling bound is asserted.
\citep{Oliveira2018Approximation,Razborov1989}
\declat{503ce2b0d3d0b5cc3dd3fa793aab840d07e448fa}{Complexity/Computability/Circuit/Fusion/Sampling.lean}{308}{Cslib.Circuits.Fusion.le_booleanPairSample_expect}
\declat{503ce2b0d3d0b5cc3dd3fa793aab840d07e448fa}{Complexity/Computability/Circuit/Fusion/Sampling.lean}{328}{Cslib.Circuits.Fusion.exists_pair_weight_ge}
\declat{503ce2b0d3d0b5cc3dd3fa793aab840d07e448fa}{Complexity/Computability/Circuit/Fusion/Sampling.lean}{355}{Cslib.Circuits.Fusion.fixedDistribution_counting_bound}
\end{result}

\begin{result}
\subsection{Linear cover complexity of random bipartite graphs}
Fix $n$ supplied subsets of a finite carrier. The number of targets admitting
an unrestricted semifilter cover with at most $t$ pairs is at most
$2^{(2t+1)(n+t+1)}$. Each cover is encoded by the inclusion flags used in
quadratic closure propagation, and decoding is proved to recover its target;
no incompressibility or encoding hypothesis remains.
For a uniform labelled graph $G\subseteq[N]\times[N]$, take all $2N$ rows and
columns as supplied sets and write $\rho(G)$ for unrestricted cover complexity.
Every graph has $\rho(G)\le N$, and
\[
\Pr[\rho(G)\le t]\le\frac{2^{(2t+1)(2N+t+1)}}{2^{N^2}}.
\]
In particular, for every $N\ge16$,
\[
\Pr[\lfloor N/16\rfloor<\rho(G)\le N]\ge1-2^{-N}.
\]
The distribution is uniform over all $2^{N^2}$ bipartite graphs, including
empty and complete graphs. This is a finite quantitative form of the
linear high-probability bound of Cavalar--Oliveira, not an explicit hard-graph
construction.\citep{CavalarOliveira2025}
\declat{5eb6d9a109234ce5579765b93bb24b7cfaca6e60}{Complexity/Computability/Circuit/Fusion/Encoding.lean}{139}{Cslib.Circuits.Fusion.card_small_covers_le}
\declat{5eb6d9a109234ce5579765b93bb24b7cfaca6e60}{Complexity/Computability/Circuit/Fusion/Graph/Random.lean}{91}{Cslib.Circuits.Fusion.Graph.pairCoverComplexity_le_card}
\declat{5eb6d9a109234ce5579765b93bb24b7cfaca6e60}{Complexity/Computability/Circuit/Fusion/Graph/Random.lean}{102}{Cslib.Circuits.Fusion.Graph.smallCoverProbability_le}
\declat{5eb6d9a109234ce5579765b93bb24b7cfaca6e60}{Complexity/Computability/Circuit/Fusion/Graph/Random.lean}{139}{Cslib.Circuits.Fusion.Graph.random_cover_complexity}
\end{result}

\begin{result}
\subsection{Semi-ultrafilter bounds for universal verifiers}
Let $g:\{0,1\}^{n+a}\to\{0,1\}$ and define
$f(x)=1$ exactly when $g(x,y)=1$ for every $a$-bit auxiliary assignment $y$.
Suppose a native monotone circuit computes $g$ from the positive and negative
literals of all $n+a$ inputs, with $t$ binary AND gates. There is a pair cover
of length exactly $t$ defeating every semi-ultrafilter above a true input of
$f$, with respect to its primary literals. Hence
\[
\rho_{\mathrm{ultra}}(f;\text{primary literals})\le t.
\]
A semi-ultrafilter accepts every set or its complement, possibly both. The bound
also holds for the minimum AND cost of implementing this fixed verifier.
OR gates, constants, supplied literals, wiring and fanout are free; arbitrary
numbers of auxiliary bits are allowed, including zero. This proves the lower-bound
direction of the conondeterministic characterization credited to Karchmer;
it does not construct a universal verifier from a cover or prove the converse.
\citep{CavalarOliveira2025}
\declat{fc111b65726f9b5bc721b9c63bef8a218fcf4b35}{Complexity/Computability/Circuit/Fusion/Conondeterministic.lean}{137}{Cslib.Circuits.Fusion.pairCoverComplexity_forallGuess_le_cost}
\end{result}
